\chapter{Literature Review}
\label{ch:literature}

\section{Introduction}

Automatic Music Transcription (AMT) has been studied for several decades, but the great majority of published work has concentrated on instruments and traditions for which large annotated datasets already exist. This chapter reviews the research that shaped the design decisions made in this thesis: general approaches to AMT, the particular difficulties associated with monophonic instruments and low-resource corpora, acoustic feature representations used for note-level analysis, deep learning architectures used for onset detection and classification, prior work on Indian classical music, and techniques for generating synthetic training data. Read together, these strands motivate the two central choices made in this work: building a synthetic melody corpus from isolated flute recordings, and pairing an RCNN-based segmentation stage with a CQT-driven CNN classifier. The chapter is organised as follows. Section~\ref{sec:amt-general} reviews general AMT frameworks; Section~\ref{sec:mono} covers monophonic transcription and low-resource instrument corpora; Sections~\ref{sec:features} and~\ref{sec:architectures} examine feature representations and neural architectures respectively; Section~\ref{sec:icm} surveys work specific to Indian classical music; Section~\ref{sec:synthetic} addresses synthetic data generation; and Section~\ref{sec:gap} states the research gap this thesis addresses.

\section{Automatic Music Transcription: General Approaches}
\label{sec:amt-general}

The task of AMT, converting an audio recording into a symbolic representation of pitch, onset and duration, has been described by Benetos et al.~\cite{ref1} as one of the long-standing open problems in music information retrieval, owing to the range of sub-tasks it subsumes: pitch estimation, onset and offset detection, instrument identification and, for polyphonic material, source separation. A later overview by the same group~\cite{ref2} organises the field into frame-level, note-level, stream-level and notation-level transcription, a hierarchy that is useful here because the present work operates at the note level: it must recover boundaries and labels, not merely a frame-wise pitch curve. A recent systematic survey by Jamshidi et al.~\cite{ref3} traces how the field moved from non-negative matrix factorization methods toward neural approaches, and notes that convolutional and recurrent architectures now dominate published AMT systems because they can be trained end to end on spectrogram-like inputs without hand-engineered pitch-tracking rules.

Piano transcription has received the most attention within AMT, largely because datasets such as MAPS and MAESTRO provide thousands of hours of aligned audio-MIDI pairs. Hawthorne et al.~\cite{ref4} introduced the Onsets and Frames architecture, which pairs a convolutional-recurrent onset detector with a framewise pitch detector and conditions the latter on the former's predictions, an approach that substantially improved note F1-scores over prior systems on the MAPS benchmark. The core idea, that separating onset detection from sustained-pitch classification and then fusing the two produces cleaner note boundaries than either sub-task alone, carries over to the two-stage segmentation-then-classification design adopted in this thesis, even though the underlying instrument, dataset scale and network types differ substantially. The transfer is not automatic, however: the onset head in Onsets and Frames exploits the sharp percussive attack of a struck string, and the flute produces no such attack. Section~\ref{sec:architectures} returns to this point.

\section{Monophonic Transcription and Low-Resource Instrument Corpora}
\label{sec:mono}

Before neural methods became standard, monophonic pitch estimation was handled by signal-processing algorithms that remain strong baselines. YIN~\cite{ref5} estimates the fundamental period from the autocorrelation-derived difference function in the time domain, while SWIPE~\cite{ref6} matches the spectrum against a sawtooth-inspired kernel in the frequency domain; pYIN~\cite{ref7} later added a hidden Markov layer over YIN candidates to produce smoother pitch tracks. These methods are accurate on clean sustained tones and require no training data, but they estimate pitch rather than notes: converting a frame-wise pitch curve into discrete note events still needs heuristic segmentation, which degrades under rapid articulation, ornamentation and expressive pitch movement. That gap between pitch tracking and note segmentation is one of the reasons this thesis treats segmentation as a learned task rather than a post-processing step.

Much of the AMT literature targets polyphonic material, where several simultaneous pitches have to be disentangled. Monophonic transcription is simpler in principle, since only one note sounds at a time, but it remains far from solved for instruments that lack large labeled corpora. Sofronievski and Gerazov~\cite{ref8} built Scorpiano, a signal-processing pipeline for monophonic piano transcription that combines onset detection, tempo estimation and pitch tracking, and reported accuracy competitive with commercial neural systems while requiring no training data at all, a useful reminder that classical DSP techniques remain a viable baseline when annotated audio is scarce. Th\'{o}me and Ahlb\"{a}ck~\cite{ref9} took a different route, training a convolutional-recurrent (ConvLSTM) pitch detector without an explicit musical language model and reporting an F-measure of 83\% on the 2007 MIREX multi-F0 development set for a woodwind ensemble recording containing bassoon, clarinet, flute, horn and oboe. Their system is polyphonic rather than monophonic, so it is cited here not as a monophonic baseline but as evidence that a learned, timbre-agnostic representation handles flute-family material without instrument-specific tuning. The same argument is made more directly by Bittner et al.~\cite{ref10}, whose lightweight instrument-agnostic model is trained jointly for onsets, multipitch and note activations and is explicitly motivated by low-resource, instrument-agnostic transcription; their results show that a single compact model built on a harmonic CQT front end generalises across instrument families at a small cost in accuracy relative to specialised systems.

Two studies deal directly with monophonic transcription of underrepresented instruments and stand as the closest precedents to the present work. Loresco et al.~\cite{ref11} built a deep learning pipeline to transcribe the diwdiw-as, a Filipino bamboo panpipe, using a dataset collected specifically for the study, and reported 96.11\% pitch accuracy and 94.07\% note-duration accuracy. Their result shows that competitive transcription accuracy is attainable for a folk monophonic instrument without access to a MAPS-scale corpus, provided the pipeline is built around the available data rather than borrowed wholesale from piano AMT. Jancovic, K\"{o}k\"{u}er and Baptiste~\cite{ref12}, working on ornamented Irish traditional flute recordings, used hidden Markov models over MFCC features, with separate models for notes and for single-note ornaments, and reported that the soft onsets and very short durations of ornaments remain a major source of transcription error even when the underlying pitch content is estimated well. This is the most directly relevant caution for the present study: the instrument family is the same, and the synthetic melody generator used here introduces comparable articulation variation through its transition styles and DSP augmentations. Taken together, these papers frame the problem this thesis addresses: transcription systems for culturally significant but data-poor monophonic instruments cannot simply reuse piano-scale training regimes and must be designed around the data constraints from the outset.

\section{Acoustic Feature Representations}
\label{sec:features}

Feature choice has a direct effect on segmentation and classification accuracy, and the transcription literature contains repeated comparisons of MFCC, Mel-spectrogram and Constant-Q Transform (CQT) representations. The CQT, in the efficient invertible formulation of Sch\"{o}rkhuber and Klapuri~\cite{ref13}, places frequency bins on a geometric grid, so that the spacing between adjacent semitones is constant across the whole range. Two consequences matter for note-level work. First, a change of pitch becomes a fixed translation along the frequency axis, which suits convolutional kernels that share weights across that axis; the same is not true of a linear or Mel-scaled spectrogram, where the distance between semitones shrinks in the upper register. Second, the constant-Q property gives finer frequency resolution in the low register, where the fundamental of a bamboo flute lies, at the cost of coarser time resolution there. MFCCs, by contrast, were designed for speech: the cepstral truncation that makes them compact for phonetic modelling deliberately smooths away the fine harmonic structure that carries pitch identity.

Empirical comparisons support this reasoning without making it universal. Cheuk et al.~\cite{ref14} evaluated several input representations for neural transcription systems under matched conditions and found that the choice of time-frequency representation changes transcription accuracy materially rather than marginally. Hernandez-Olivan et al.~\cite{ref15} compared a Mel-based Onsets and Frames model against a CQT-based salience model across the instrument families of the Slakh2100 dataset, which includes a flute subset, and reported that the CQT-based model degrades far less on instruments whose onsets are weaker than those of piano or plucked strings; training on isolated instrument families improved results further. Within Indian classical music specifically, Hebbar and Jagtap~\cite{ref16} compared preprocessing choices against several deep architectures for Carnatic raga recognition, establishing a precedent for treating the feature representation as a controlled experimental variable rather than a fixed design decision. The counterexample is equally instructive: Wolf-Monheim~\cite{ref17}, evaluating spectral and rhythm features for general-purpose audio classification, found Mel-scaled spectrograms and MFCCs to outperform CQT and CENS chromagrams. Taken together, these results indicate that the best representation is task-dependent, that CQT tends to win where the target is note-level pitch content on a musical scale, and that the question should therefore be settled experimentally for the flute rather than assumed. That is the approach adopted in Chapter~\ref{ch:results} of this thesis.

\section{Neural Architectures for Onset Detection and Note Segmentation}
\label{sec:architectures}

Onset detection is the sub-task most directly analogous to the segmentation stage of this thesis. Schl\"{u}ter and B\"{o}ck~\cite{ref18} showed that a CNN operating on short log-Mel spectrogram excerpts could outperform the hand-designed spectral-flux detectors that had been the state of the art for over a decade, while requiring far less manual signal-processing tuning. Their result is frequently cited as early evidence that, for onset detection specifically, learned convolutional features generalise better than engineered ones. Eyben et al.~\cite{ref19} had earlier shown that bidirectional LSTMs trained on auditory spectral features could detect onsets across a wide range of instrument timbres without retuning, establishing that recurrent modelling of temporal context, and not just local spectral shape, is useful for this task. \c{C}ak{\i}r et al.~\cite{ref20} combined both ideas in a convolutional-recurrent architecture for polyphonic sound event detection and, in follow-up work, extended it with a frame-level attention mechanism so the network could down-weight irrelevant or silent frames. This is the same architectural family, CNN feature extraction, RNN temporal modelling and attention-based frame weighting, that underlies the RCNN-with-attention segmentation model used in this thesis. Attention has since been applied inside AMT itself: Wu et al.~\cite{ref21} use self-attention to frame multi-instrument transcription as an instance segmentation problem, treating note events as objects to be delineated rather than frames to be labelled independently.

A second and less frequently cited strand is more directly applicable to the flute, because the flute produces pitched, non-percussive onsets with a slow amplitude rise rather than a sharp transient. Stasiak et al.~\cite{ref22} fused several onset detection functions through a multi-layer perceptron rather than thresholding any single function, and evaluated across seventeen instruments and ensembles. Mo\'{n}ko and Stasiak~\cite{ref23} then applied a CNN onset detector specifically to bowed string recordings, the canonical soft-onset case, and found that a network trained on mixed material transferred well to solo recordings of a sustained-tone instrument. Mounir et al.~\cite{ref24} approached the same problem from the feature side, proposing the NINOS\textsuperscript{2} spectral sparsity measure and showing that it outperforms logarithmic spectral flux precisely where conventional detectors are weakest: sustained-string instruments, repeated notes and vibrato passages. These three papers matter for the present work because they identify soft onsets, repeated notes and vibrato as the failure modes of onset-based segmentation, and all three conditions occur in the synthetic flute melodies used here, both through the slurred and crossfaded transitions and through the vibrato augmentation. None of the reviewed onset-detection papers, however, evaluate on Indian classical instruments or on a corpus built from synthetically concatenated single-note recordings, which is the specific setting this thesis targets.

\section{Note and Raga Recognition in Indian Classical Music}
\label{sec:icm}

Computational work on Indian classical music (ICM) has been reviewed in detail by Singh and Biswas~\cite{ref25}, who note that the field has concentrated on a small set of tasks, raga recognition foremost among them, and that note-level analysis remains comparatively underdeveloped. A difficulty specific to this repertoire is tonic (\'{s}ruti) normalisation. Unlike Western fixed-pitch instruments, the absolute frequency of a given swara depends on the tonic (Sa) chosen by the performer, so the same swara label maps to different absolute pitches across recordings. Bellur et al.~\cite{ref26} proposed pitch-histogram and group-delay methods for identifying the tonic, and Gulati et al.~\cite{ref27} later compared the main approaches on a common evaluation set; Singh et al.~\cite{ref28} took the alternative route of building the tonic into the representation itself through a shruti-based chromagram. Many raga-recognition systems avoid the problem entirely by working with tonic-independent chroma or pitch-interval features. The segmentation stage of this thesis instead operates on absolute CQT contours, because boundary detection depends on spectral change rather than on scale-degree identity; the consequence is that the classification stage inherits a dependence on the tonic of the source recordings, and tonic normalisation is treated in Chapter~\ref{ch:limitations} as a limitation and a direction for future work rather than as a solved problem.

At the level of melodic modelling, early work by Dighe et al.~\cite{ref29} used chromagram-derived swara histograms with classical classifiers, showing that swara occurrence and transition statistics carry strong discriminative information even without deep models. More recent systems use neural sequence models: DeepSRGM~\cite{ref30} applies an LSTM-based recurrent network to Carnatic recordings for raga classification and ranking, PhonoNet~\cite{ref31} uses a multi-stage CNN pipeline for Hindustani raga identification, and Shah et al.~\cite{ref32} report high accuracy on a ten-raga subset of the CompMusic collection using CNNs over raw spectrograms. Singh and Arora~\cite{ref33} revisit the task with explainability in mind and show that the regions their model attends to correspond to short melodic phrases rather than isolated notes, a useful caution for any system, including the one proposed here, that treats swara boundaries as if they were as sharply defined as piano note onsets. Two studies by Nuttall et al.~\cite{ref34,ref35} work at the note level directly, modelling how a given svara's acoustic realisation is shaped by its melodic context and by gamaka ornamentation in Carnatic music, through deep clustering and embedding methods respectively. Their finding that a svara's spectral shape is not fixed but depends heavily on its neighbours is directly relevant to the classification stage of this thesis, since the synthetic melody generator introduces comparable context-dependent variation through its transition-style and DSP augmentation parameters.

The closest published system to the present work is that of Singh et al.~\cite{ref36}, who apply a convolutional-recurrent network with a CTC objective to transcription of Indian classical music. Their use of CTC is significant: it removes the need for frame-aligned labels, but in exchange the model predicts a sequence of symbols without committing to explicit onset and offset times. The framework proposed in this thesis takes the opposite position, predicting frame-level boundaries directly, which is what makes note-duration and timing information recoverable. Neither their work nor the raga literature surveyed above addresses the bamboo flute specifically, and none of it evaluates frame-level segmentation of continuous melodies into individual swaras.

\section{Addressing Data Scarcity through Synthetic Data Generation}
\label{sec:synthetic}

Because annotated recordings of Indian classical flute performances are scarce, the synthetic melody generation strategy adopted in this thesis sits within a broader body of work on synthetic and recombined training data. The closest precedent within music transcription is the Slakh dataset~\cite{ref37}, which renders MIDI transcriptions through virtual instruments to obtain large volumes of audio with exact, automatically derived annotations, and which has since become a standard testbed for multi-instrument transcription. Slakh establishes the central argument for procedurally generated audio in this field: when ground truth is produced by the generator rather than by a human annotator, it is exact by construction and available at any scale. Its principal weakness is equally well known, in that the timbre of a sampled virtual instrument is not the timbre of a live performance.

The approach used here inverts that trade-off. Rather than synthesising audio from symbolic input, isolated real flute recordings are concatenated into continuous melodies with programmatically controlled transitions and DSP augmentations, so that every acoustic unit in the corpus is a real recording of the target instrument while the boundary annotations remain exact. Work elsewhere in the audio community follows the generative rather than the recombinative paradigm: Ghosh et al.~\cite{ref38} augment small audio classification datasets with a text-to-audio diffusion model whose outputs are aligned to the target distribution through preference optimisation, and Ibaraki and Chiang~\cite{ref39} generate training data for low-resource speech recognition by pairing novel text with text-to-speech synthesis. Both report that carefully constructed synthetic data substitutes for scarce real recordings rather than merely padding dataset size, which supports the general premise, but neither preserves the acoustic identity of a specific physical instrument in the way that recombination does. The limitation of the recombinative approach should be stated plainly, since it bears on how the results of this thesis are interpreted: concatenated melodies reproduce the timbre of the source notes but not the continuous breath control, phrase-level dynamics and gamaka-like pitch movement of a live performance, so accuracy measured on this corpus is an upper bound on what would be obtained on unconstrained recordings. Chapter~\ref{ch:limitations} discusses this gap and the evaluation on held-out real audio used to bound it.

\section{Summary and Research Gap}
\label{sec:gap}

Across this body of work, three patterns stand out. First, deep learning based AMT is well developed for piano and, increasingly, for polyphonic multi-instrument settings, but monophonic transcription for non-Western or folk instruments remains scattered across a handful of isolated studies, most of which build small custom datasets in a manner similar to the approach taken here. Second, CNN, RNN and attention architectures originally developed for onset detection and polyphonic sound event detection have been applied to soft-onset instruments such as bowed strings, but not, to the best of our knowledge, to frame-level segmentation of continuous Indian classical flute melodies into swara boundaries, even though the underlying ideas of frame-level attention, CQT input and CNN-RNN fusion transfer naturally to this setting. Third, the Indian classical music literature has concentrated on raga-level classification and, more recently, on context-dependent svara realisation and CTC-based sequence transcription, but has not combined a segmentation-first pipeline with a systematic feature comparison and an evaluation on synthetically generated continuous melodies.

The framework developed in the remainder of this thesis sits at the intersection of these three gaps. Specifically, it makes three contributions:

\begin{enumerate}
    \item[(i)] it adapts attention-based convolutional-recurrent segmentation, established in the general AMT and sound event detection literature, to the understudied domain of the monophonic Indian classical flute, where onsets are soft and note boundaries are not marked by percussive transients;
    \item[(ii)] it evaluates CQT systematically against Mel-spectrogram, MFCC, pitch and RMS energy features for this specific timbre and task, rather than adopting a representation by convention; and
    \item[(iii)] it introduces a recombination-based synthetic melody generator that removes the manual annotation bottleneck while keeping every acoustic unit grounded in a real flute recording, and that is instrument-independent by construction.
\end{enumerate}
